\subsubsection{Statement and Proof of the Theorem}\label{subsec:ck-statement-proof}

Consider the system of partial differential equations
\begin{equation}\label{eq:ck-quasilinear-system}
\dfrac{\partial u_i}{\partial x_n}=\sum_{k=1}^{n-1}\sum_{j=1}^Na_{ij}^k(\bm{p})\dfrac{\partial u_j}{\partial x_k}+b_i(\bm{p}),\qquad i=1,\ldots,N,
\end{equation}
where $\bm{p}$ stands for the vector $(x_1,\ldots,x_{n-1},u_1,\ldots,u_N)$, with the initial conditions
\begin{equation}\label{eq:ck-zero-initial-data}
u_i=0\quad\text{on }x_n=0,\qquad i=1,\ldots,N.
\end{equation}
We shall establish the following local existence result.

\begin{Theorem}
\label{thm:cauchy-kovalevskaya-local-existence}
Let the functions $a_{ij}^k$ and $b_i$ be real analytic at the origin of $\mathbb{R}^{N+n-1}$. Then the system \eqref{eq:ck-quasilinear-system} with initial conditions \eqref{eq:ck-zero-initial-data} has a unique (among real analytic functions) system of solutions $u_i$ that is real analytic at the origin.
\end{Theorem}

\begin{Proof}We begin by formally computing all derivatives of the $u_i$ at the origin. From \eqref{eq:ck-zero-initial-data} it follows that all tangential derivatives of all orders are zero. Hence the only nonzero first derivative is in the $x_n$ direction, and \eqref{eq:ck-quasilinear-system} leads to $\partial u_i/\partial x_n(\bm{0})=b_i(\bm{0})$. Next, we can differentiate \eqref{eq:ck-quasilinear-system} to obtain second derivatives of $u_i$. We find
\begin{equation}\label{eq:ck-second-derivative-recurrence}
\begin{aligned}
\dfrac{\partial^2u_i}{\partial x_n\partial x_l}(\bm{0})
&=\sum_{k=1}^{n-1}\sum_{j=1}^Na_{ij}^k(\bm{0})\dfrac{\partial^2u_j}{\partial x_k\partial x_l}(\bm{0})\\
&\quad+\sum_{j=1}^{N+n-1}\dfrac{\partial b_i}{\partial p_j}(\bm{0})\dfrac{\partial p_j}{\partial x_l}(\bm{0}),
\end{aligned}
\end{equation}
which yields, for $l=1,\ldots,n-1$,
\begin{equation}\label{eq:ck-tangential-second-derivatives}
\dfrac{\partial^2u_i}{\partial x_n\partial x_l}=\dfrac{\partial b_i}{\partial p_l},
\end{equation}
and, for $l=n$, $i=1,\ldots,N$,
\begin{equation}\label{eq:ck-second-normal-derivative}
\begin{aligned}
\dfrac{\partial^2u_i}{\partial x_n^2}
&=\sum_{k=1}^{n-1}\sum_{j=1}^Na_{ij}^k(\bm{0})\dfrac{\partial^2u_j}{\partial x_k\partial x_n}(\bm{0})
+\sum_{j=1}^N\dfrac{\partial b_i}{\partial p_{n-1+j}}(\bm{0})\dfrac{\partial u_j}{\partial x_n}(\bm{0})\\
&=\sum_{k=1}^{n-1}\sum_{j=1}^Na_{ij}^k(\bm{0})\dfrac{\partial b_j}{\partial p_k}(\bm{0})
+\sum_{j=1}^N\dfrac{\partial b_i}{\partial p_{n-1+j}}(\bm{0})b_j(\bm{0}).
\end{aligned}
\end{equation}


Proceeding in a similar fashion, we can compute derivatives of all orders. The resulting expression for any derivative of $u_i$ at the origin is a polynomial with positive coefficients involving derivatives of $a_{ij}^k$ and $b_i$ at the origin.

We can thus construct a formal Taylor series for the $u_i$. To show that this Taylor series converges, we use the method of majorants. Let
\begin{equation}\label{eq:ck-coefficient-majorants}
a_{ij}^k(\bm{p})\preceq A_{ij}^k(\bm{p}),\qquad b_i(\bm{p})\preceq B_i(\bm{p}),
\end{equation}
and let $U_i$ be the solution of the problem
\begin{equation}\label{eq:ck-majorant-system}
\dfrac{\partial U_i}{\partial x_n}=\sum_{k=1}^{n-1}\sum_{j=1}^NA_{ij}^k(\bm{p})\dfrac{\partial U_j}{\partial x_k}+B_i(\bm{p}),\qquad i=1,\ldots,N,
\end{equation}
with the initial conditions
\begin{equation}\label{eq:ck-majorant-zero-data}
U_i=0\quad\text{on }x_n=0,\qquad i=1,\ldots,N.
\end{equation}
Then clearly $|D^\alpha u_i(\bm{0})|\leq D^\alpha U_i(\bm{0})$ for any $\alpha$; hence the formal power series for $u_i$ converges if $U_i$ is analytic.

It remains to construct appropriate functions $A_{ij}^k$ and $B_i$. Let us assume that $a_{ij}^k$ and $b_i$ are in $C_{M,r}(\bm{0})$. Then we can find a majorant from \eqref{eq:analytic-rational-majorant}, i.e.,
\begin{equation}\label{eq:ck-rational-coefficient-majorants}
A_{ij}^k=B_i=\dfrac{Mr}{r-x_1-\cdots-x_{n-1}-U_1-\cdots-U_N}.
\end{equation}
That is, we have to consider the initial-value problem
\begin{equation}\label{eq:ck-rational-majorant-system}
\dfrac{\partial U_i}{\partial x_n}=\dfrac{Mr}{r-x_1-\cdots-x_{n-1}-U_1-\cdots-U_N}\left(1+\sum_{k=1}^{n-1}\sum_{j=1}^N\dfrac{\partial U_j}{\partial x_k}\right),
\end{equation}
with initial conditions \eqref{eq:ck-majorant-zero-data}.

A solution of this problem can be found in the form $U_i(\bm{x})=V(x_1+\cdots+x_{n-1},x_n)$; the resulting equation for $V(s,t)$ is
\begin{equation}\label{eq:ck-scalar-majorant-problem}
V_t=\dfrac{Mr}{r-s-NV}(1+N(n-1)V_s),\qquad V(s,0)=0.
\end{equation}


This equation can be solved explicitly. We first note that the characteristic lines $s=s(t)$ are given by the equation
\begin{equation}\label{eq:ck-majorant-characteristics}
\dfrac{\mathrm{d}s}{\mathrm{d}t}=-\dfrac{(n-1)NMr}{r-s-NV(s(t),t)}.
\end{equation}


Along a characteristic, we have
\begin{equation}\label{eq:ck-majorant-along-characteristic}
\dfrac{\mathrm{d}}{\mathrm{d}t}V(s(t),t)=\dfrac{Mr}{r-s(t)-NV(s(t),t)}=-\dfrac1{(n-1)N}\dfrac{\mathrm{d}s}{\mathrm{d}t}.
\end{equation}
Integration with respect to $t$ yields
\begin{equation}\label{eq:ck-majorant-characteristic-integral}
V(s(t),t)=-\dfrac1{(n-1)N}s(t)+\dfrac1{(n-1)N}s(0).
\end{equation}


By inserting \eqref{eq:ck-majorant-characteristic-integral} into \eqref{eq:ck-majorant-characteristics}, we find
\begin{equation}\label{eq:ck-reduced-characteristic-equation}
\dfrac{\mathrm{d}s}{\mathrm{d}t}=-\dfrac{(n-1)MNr}{r-s+\dfrac{s}{n-1}-\dfrac{s(0)}{n-1}},
\end{equation}
which can be integrated to obtain
\begin{equation}\label{eq:ck-initial-characteristic-coordinate}
s(0)=\dfrac{n-1}{n}\left[r+\dfrac1{n-1}s-\sqrt{(r-s)^2-2nNMrt}\right].
\end{equation}
We insert this into \eqref{eq:ck-majorant-characteristic-integral} and finally find
\begin{equation}\label{eq:ck-explicit-scalar-majorant}
V(s,t)=\dfrac1{Nn}\left(r-s-\sqrt{(r-s)^2-2nNMrt}\right).
\end{equation}
This expression is analytic at $s=t=0$. This concludes the proof.\end{Proof}

The theorem guarantees the existence of a solution in a neighborhood of each point on the initial surface. By taking the union of all these neighborhoods, we obtain the existence of a solution in an open set containing the initial surface. We also note that the Taylor series for the solution is guaranteed to converge in a ball whose radius depends only on $n$, $N$, $M$, and $r$.

We emphasize that the theorem is strictly local in character, that is, the solution may cease to exist or at least lose analyticity at some finite value of $x_n$. We shall see examples of this later when we study development of shocks. Also, the theorem does not rule out the existence of nonanalytic solutions even in a neighborhood of the initial surface. Uniqueness is only guaranteed within the class of analytic functions. Holmgren's theorem, proved in the next section, asserts that for linear equations the analytic solution is unique in a larger function class.
